% Generated from locale/tr/content/set-theory/z/story.tex; do not hand-edit.


\olfileid[tr]{sth}{z}{story}
\olsection{Anlatının Daha Ayrıntılı Biçimi}

\olref[story][approach]{sec} içinde Schoenfield'ın küme oluşturma süreci
betimlemesini aktardık. Şimdi bu anlatıyı biraz daha kesinleştirmek için
birkaç ilke daha yazmak istiyoruz. İlkeler şunlardır:
\begin{enumerate}
\item[] \stageshier. Her küme bir evrede oluşturulur.
\item[] \stagesord. Evreler sıralıdır: Bazıları diğerlerinden \emph{önce}
gelir.\footnote{Aslında evrelerin \emph{iyi sıralı} olduğunu örtük olarak
varsayacağız. Bunun ne anlama geldiği \olref[ordinals][]{chap} içinde
açıklanır. Bu önemli bir varsayımdır. Hatta \citet{Scott1974}'e ait çok
zekice bir teknik kullanılarak bu varsayımdan \emph{kaçınılabilir}, ardından
varsayım \emph{türetilebilir}. (Bu aynı zamanda neden bir başlangıç evresi
olduğunu düşünmemiz gerektiğini de açıklar.) Burada buna giremeyiz; daha
fazlası için bkz. \citet{ButtonLT1}.}
\item[] \stagesacc. Her $S$ evresi ve $S$ evresinden \emph{önce}
oluşturulmuş gelişigüzel kümeler için: $S$ evresinde, elemanları tam olarak
bu kümeler olan bir küme oluşturulur. $S$ evresinde başka hiçbir şey
oluşturulmaz.
\end{enumerate}
Bunlar biçimsel olmayan ilkelerdir; fakat Zermelo küme teorisinin çeşitli
aksiyomlarını haklı çıkarmak için bunları kullanabileceğiz.

(Bir uyarı eklemeliyiz. Tümüyle standart ve standart adlara sahip bazı
aksiyomlar sunacak olsak da az önce sunduğumuz italik ilkelerin yazında özel
bir adı yoktur. Yalnızca yararlı olacağını umduğumuz adlar verdik.)

